% ============================================================================
% TECHNICAL APPENDIX 5A: The Mathematics and Psychology of Optimal Human–AI Asks
% Formal foundations for intervention timing, frequency, and format design
% ============================================================================
% Source: /chapters/ch5/Ch5A_final.md
% ============================================================================

\chapter*{Technical Appendix 5A: Optimal Ask Design}
\addcontentsline{toc}{chapter}{Technical Appendix 5A: Optimal Ask Design}

\begin{chapterquote}{Formal foundations for Chapter 5}
Chapter~5 introduced three design levers for human-in-the-loop asks: timing, frequency, and format. This appendix provides mathematical frameworks for optimising each, drawing on information theory, cognitive psychology, and reinforcement learning.
\end{chapterquote}

% ============================================================================
\section{Information-Theoretic Foundations of Ask Design}
% ============================================================================

\subsection{Information Value of a Human Response}

The value of asking a human is the expected reduction in uncertainty about the correct action. Let $Y$ be the correct decision (unknown to the system), $H$ the human's response, and $X$ the system's current information. The \term{information gain} from asking is:

\[
\text{IG}(H;\, Y \mid X) = \mathbb{H}(Y \mid X) - \mathbb{H}(Y \mid H, X)
\]

where $\mathbb{H}(\cdot)$ denotes entropy. A high information gain means the human's response substantially reduces uncertainty about the right action.

\paragraph{Optimal ask criterion.} Ask a human when:

\[
\text{IG}(H;\, Y \mid X) > \frac{C_\text{ask}}{V_\text{decision}}
\]

where $C_\text{ask}$ is the cost (time, attention, trust) and $V_\text{decision}$ is the value of making the correct decision.

\subsection{Query Optimisation}

The \term{optimal query} $Q^*$ maximises information gain subject to cognitive load constraints:

\[
Q^* = \arg\max_{Q \in \mathcal{Q}} \bigl[\text{IG}(H_Q;\, Y \mid X) - \lambda \cdot \text{CogLoad}(Q)\bigr]
\]

where $\text{CogLoad}(Q)$ is the cognitive effort required to answer query $Q$ and $\lambda$ is a cost--effort tradeoff parameter.

\begin{examplebox}[Simplicity vs.\ Information]
A simpler query (lower cognitive load) is only preferable if its information loss is proportionally small. The 911~example from Chapter~5 illustrates this: the revised question had \emph{lower} cognitive load and \emph{higher} information gain simultaneously---a rare but achievable combination.
\end{examplebox}

\subsection{Human--AI Complementarity}

The value of the full human--AI system depends on the \emph{complementarity} between human and AI knowledge:

\[
I(H_\text{human};\, Y \mid H_\text{AI}, X) = \mathbb{H}(Y \mid H_\text{AI}, X) - \mathbb{H}(Y \mid H_\text{human}, H_\text{AI}, X)
\]

This is the additional uncertainty reduction from the human, given the AI has already processed $X$. High complementarity means the human knows things the AI does not---this is when HITL intervention is most valuable.

% ============================================================================
\section{Alert Fatigue: A Dynamic Model}
% ============================================================================

\subsection{Trust as a State Variable}

Let $\tau_t \in [0, 1]$ represent the human's trust in the system's alerts at time~$t$. Trust evolves dynamically:

\[
\tau_{t+1} = \tau_t
  + \alpha \cdot \mathbf{1}[\text{alert}_t \text{ was valid}]
  - \beta  \cdot \mathbf{1}[\text{alert}_t \text{ was false}]
\]

where $\alpha$ is the trust gain from a true alert and $\beta$ is the trust loss from a false alarm. Empirically $\beta > \alpha$: trust erodes faster than it builds \autocite{lee2004trust}. Effective alert value therefore decays with sustained false-positive rates:

\[
V_\text{effective}(t) = \tau_t \cdot V_\text{intrinsic}
\]

\subsection{Precision--Fatigue Tradeoff}

Define $r$ as the alert rate, $p(r)$ as precision (decreasing in $r$), $\tau(r)$ as equilibrium trust (also decreasing in $r$), and $\rho(\tau)$ as response probability. The effective information extracted per hour:

\[
I_\text{effective}(r) = r \cdot p(r) \cdot \tau(r) \cdot \rho\bigl(\tau(r)\bigr) \cdot V_\text{per alert}
\]

\begin{lstlisting}[style=python, caption={Effective information vs.\ alert rate}]
from scipy.optimize import minimize_scalar

def effective_information(r, precision_decay=0.1,
                           trust_decay=0.05, base_value=1.0):
    precision = 1 / (1 + precision_decay * r)
    false_alarm_rate = r * (1 - precision)
    trust = 1 / (1 + trust_decay * false_alarm_rate)
    response_prob = trust ** 0.5
    return r * precision * trust * response_prob * base_value

result = minimize_scalar(
    lambda r: -effective_information(r),
    bounds=(0.1, 100), method='bounded'
)
print(f"Optimal alert rate: {result.x:.2f} alerts/hour")
\end{lstlisting}

\subsection{Habituation and Recovery}

Consecutive alerts induce faster degradation:

\[
\tau_t = \tau_{\text{baseline}} \cdot e^{-\gamma \cdot n_\text{recent}}
\]

where $n_\text{recent}$ is the number of alerts in the recent window. Recovery is slower:

\[
\tau_{t+T} = \tau_t + (1 - \tau_t)\bigl(1 - e^{-\delta T}\bigr)
\]

\begin{cautionbox}[Design Implication]
Spacing alerts recovers attention; clustering alerts degrades it faster than average-rate models predict. A well-designed HITL interface enforces a minimum inter-alert gap, not just a maximum average rate.
\end{cautionbox}

% ============================================================================
\section{Optimal Timing Models}
% ============================================================================

\subsection{Attention Windows}

Let $a(t) \in [0,1]$ be the human's attention availability. The effective response quality for an alert sent at time~$t$:

\[
Q_\text{response}(t) = Q_\text{max} \cdot a(t) \cdot \tau(t)
\]

Optimal timing: $t^* = \arg\max_t Q_\text{response}(t)$. Interruption at natural task breakpoints consistently produces higher response quality than mid-task interruption.

\subsection{The Window of Relevance Problem}

Context availability decays after an event:

\[
C(\Delta t) = e^{-\kappa \Delta t}
\]

for decay rate $\kappa$ depending on event type. Optimal alert delay balances context availability against attention window:

\[
t^* = \arg\max_t \, a(t + t_\text{event}) \cdot C(t)
\]

For financial transactions $\kappa$ is large, so $t^* \approx 0$. Asking about a charge 24~hours later yields near-chance accuracy.

% ============================================================================
\section{Format Design: Cognitive Load and Response Accuracy}
% ============================================================================

\subsection{Cognitive Load Model}

Following Sweller's cognitive load theory, response accuracy degrades beyond working memory capacity $W$:

\[
A(c, W) = \begin{cases}
  A_\text{max} & \text{if } c \leq W \\
  A_\text{max} \cdot (W/c)^\alpha & \text{if } c > W
\end{cases}
\]

\subsection{Optimal Query Granularity}

A $k$-choice query captures $\log_2 k$ bits at $O(k)$ cognitive cost:

\[
k^* = \arg\max_k \frac{\log_2 k}{k^\beta \cdot C_\text{per-choice}}
\]

For most practical settings $k^* \in [2,5]$. Binary questions are easiest but information-sparse; 4--5 choices represent the practical optimum.

% ============================================================================
\section{Feedback Integration: Closing the Learning Loop}
% ============================================================================

\subsection{Active Learning Formulation}

Each human response is a labeled data point. Select the next query to maximise expected model improvement:

\[
Q^*_\text{next} = \arg\max_Q \mathbb{E}\bigl[\Delta\text{ModelPerf} \mid H_Q,\, \mathcal{D}_\text{current}\bigr]
\]

In practice this is approximated by querying cases near the decision boundary, in underrepresented categories, or where human and model systematically disagree.

\subsection{RL Formulation of the Full HITL Loop}

The ask--learn loop formalised as an MDP \autocite{christiano2017deep}:

\begin{itemize}
    \item \textbf{State} $s_t$: model weights, trust $\tau_t$, queue depth, context
    \item \textbf{Action} $a_t$: ask with format $Q$, or act autonomously
    \item \textbf{Reward}: $r_t = \Delta\text{Acc}(a_t) - C_\text{ask}\cdot\mathbf{1}[a_t=\text{ask}] - \lambda\cdot\Delta\tau_t$
\end{itemize}

Optimal policy $\pi^*$ maximises $\mathbb{E}[\sum_t \gamma^t r_t]$, capturing the tension between short-run accuracy and long-run trust.

% ============================================================================
\section{Exercises}
% ============================================================================

\begin{Exercise}[Information Value Calculation]
A spam filter has confidence 0.65 on an email. A human reviewer has 90\% accuracy on boundary cases.
\begin{enumerate}
    \item Calculate the expected information gain from asking the reviewer.
    \item If the cost of asking is 30~seconds of reviewer time (valued at \$0.01/second), at what minimum email value must human review be worthwhile?
    \item How does information gain change if filter confidence was 0.95?
\end{enumerate}
\end{Exercise}

\begin{Exercise}[Alert Fatigue Simulation]
Using the dynamic trust model:
\begin{lstlisting}[style=python, caption={Alert fatigue simulator}]
import numpy as np

def simulate_alert_fatigue(n_alerts=300, precision=0.7,
                            alpha=0.05, beta=0.15):
    trust, history = 1.0, [1.0]
    for _ in range(n_alerts):
        is_valid = np.random.random() < precision
        trust = min(1.0, trust + alpha) if is_valid \
                else max(0.0, trust - beta)
        history.append(trust)
    return history
\end{lstlisting}
Simulate: (a)~10/hour at 70\% precision; (b)~3/hour at 90\% precision; (c)~20/hour at 90\% precision. Compare equilibrium trust levels.
\end{Exercise}

\begin{Exercise}[Optimal Query Granularity]
For a content moderation interface with 5 labels vs.\ 2 labels:
\begin{enumerate}
    \item Calculate maximum information per response for each.
    \item Assuming $\beta = 1.5$, compute $k^*$.
    \item Design an experiment measuring actual information captured using inter-annotator agreement as a proxy.
\end{enumerate}
\end{Exercise}

\begin{Exercise}[HITL Policy Optimisation]
Implement a Q-learning agent learning when to ask. State: confidence (binned 0--10), trust (high/medium/low). Actions: auto-decide, binary query, detailed query. Reward: $+10$ correct, $-1$ ask cost, $-0.5\times(1-\text{trust})$ incorrect auto-decision. Simulate 10{,}000 episodes. At what confidence level does the optimal policy switch from asking to auto-deciding?
\end{Exercise}

% ============================================================================
\section{Recommended Reading}
% ============================================================================

\subsection*{Alert Design}
\begin{itemize}
    \item \textcite{eemua191} --- Engineering Equipment and Materials Users' Association: alarm system management guidelines.
    \item \textcite{drew2013alarm} --- Attention and awareness in simulated radiological search. \emph{Psychological Science}.
\end{itemize}

\subsection*{Active Learning and HITL Feedback}
\begin{itemize}
    \item \textcite{settles2009active} --- Active learning literature survey.
    \item \textcite{christiano2017deep} --- Deep reinforcement learning from human preferences. \emph{NeurIPS}.
\end{itemize}

\bigskip
\begin{center}
\rule{0.5\textwidth}{0.4pt}
\end{center}
\bigskip

\noindent\textit{This appendix provides the formal foundations for ask design, alert fatigue, timing, format, and feedback integration as introduced in Chapter~5.}
